\section{The instrument}\label{sec:instrument}

Every index in this study is built, stored and queried by one engine, so that a difference between two variants measures the variants and not two implementations. This section describes what the engine fixes for every variant (Section~\ref{sec:engine}), how an index is placed in memory and on a storage device (Section~\ref{sec:placement}), what a variant may change, nesting included (Section~\ref{sec:designspace}), and the per-cloud search that supplies the best structure each cloud is judged against (Section~\ref{sec:search}).

\subsection{One engine for every structure}\label{sec:engine}

\paragraph{Node pool.} A built index is one array of 32-byte \emph{hot} nodes. A hot node holds the tight bounding box of the points beneath it as six floats, the index of its first child (or, for a leaf, of its first point), and one word that packs a child count with a flag telling leaves from internal nodes. Two hot nodes fill a 64-byte cache line. The children of a node are always written as one contiguous \emph{sibling group} that starts on a cache-line boundary. Everything a query never reads (the split cell a node came from, its subtree extent, its depth) lives in a parallel array of \emph{cold} nodes.

\paragraph{Points.} The points are reordered into tree order and stored as three separate coordinate arrays plus an index map back to the input, so that every subtree is one contiguous run of the arrays and a leaf scan streams through memory. Bounding boxes are tight, empty children are dropped, and chains of single children are collapsed.

\paragraph{Queries.} The engine answers axis-aligned range, count, radius and $k$-nearest-neighbour queries. Range, count and radius queries walk the tree with an explicit stack: when a node is entered, the boxes of all its children are tested while the sibling group is in cache, and the survivors are pushed; a node entirely inside the query takes its whole subtree run. Leaf scans use AVX2, eight points per step, with the same arithmetic as the scalar loop so that both return identical sets. For $k \leq 64$ the nearest-neighbour search is depth first, nearest child first; above 64 it is best first. Only the traversal is timed, and every timed pass recomputes an order-independent checksum of the returned identifiers and compares it with the first pass.

\paragraph{Node order.} Two orders of the node pool are compared (Fig.~\ref{fig:grammar}b). \emph{Depth-first}: after a node's sibling group come the subtrees of its children, one after another, so a subtree occupies one contiguous slot range. \emph{Page-blocked}: sibling groups stay contiguous, but are packed into 4\,KiB pages of 128 hot nodes breadth first from a page's first group, so that a query descending several levels stays in one page; subtrees too small to fill a page share one, and levels whose nominal fan-out exceeds a quarter page (27- and 64-way grids) keep the depth-first order, because one of their groups is already a page access. On the clouds of this study the page-blocked order reduces the distinct pages a query touches by 10 to 22\% for binary and octree hierarchies; Sections~\ref{sec:ram} and~\ref{sec:ooc} measure what that is worth in memory and on a device.

\paragraph{Pages and prefetching.} With \emph{huge pages}, every node pool and point array of four megabytes or more is aligned to two megabytes and advised to the kernel as a transparent-huge-page candidate. With \emph{child prefetch}, each child pushed on the stack has its own sibling group, or the first line of its point arrays, prefetched before the traversal moves on. Each combination of these options compiles to its own traversal routine, so the default path is unchanged by their presence.

\paragraph{Telemetry.} An instrumented pass counts, per query, the nodes entered, boxes tested, points scanned and points returned, and records the time-stamp-counter ticks between consecutive entered nodes, charged to the depth of the earlier one; the timed passes compile all of this out.

\subsection{Placement: memory, storage device, cache}\label{sec:placement}

The same index can be queried from memory or from a file on a storage device. The \emph{leaf store} writes the hot pool and the tree-ordered points of a built index into one file, every region aligned to 4\,KiB pages; a leaf's points are one contiguous run of 16 bytes per point, read back as the pages that cover it. Reads bypass the operating system's file cache (direct I/O on Linux, unbuffered I/O on Windows) and are synchronous, one at a time, so that the timed read is the device's and not a copy of memory. Three placements are measured:
\begin{itemize}[leftmargin=*, itemsep=1pt]
\item \emph{in memory}: the whole index in RAM, the reference;
\item \emph{leaves on the device}: the hierarchy in RAM, every leaf a query scans read from the device;
\item \emph{hierarchy and leaves on the device}: in addition, every sibling group a query tests is read as the page that holds it.
\end{itemize}
An optional least-recently-used cache, sized as a fraction of the store, holds leaves and hierarchy pages between reads; it is emptied before every timed pass. Traversal decisions are taken on the in-memory pool in every placement, so all placements return identical answers; the engine checks this for every query.

\subsection{The design space, and what nesting is}\label{sec:designspace}

% Figure 2: the schema grammar (left) and the pool every schema compiles into (right).
\begin{figure*}[tb]
\centering
\resizebox{\textwidth}{!}{%
\begin{tikzpicture}[x=1mm, y=1mm, font=\scriptsize]
  % ------------------------------------------------------------ left: three schemas
  \node[mdslabel, anchor=west] at (-2, 22) {(a) schemas: blocks of (primitive, split rule, levels) and one leaf capacity};
  % single
  \node[mdslabel] at (10, 16) {single};
  \node[mdsbox, minimum width=30mm, minimum height=12mm] (s1) at (10, 6) {kd-tree\\sliding midpoint\\48 levels};
  \node[mdsbox, minimum width=30mm, fill=white] (sl) at (10, -6) {leaf 128};
  \draw[mdsflow] (s1) -- (sl);
  \node[mdslabel, anchor=west] at (26, 6) {\schema{kd48l128sm}};
  % nested
  \node[mdslabel] at (58, 16) {nested};
  \node[mdsbox, minimum width=30mm] (n1) at (58, 9) {quadtree, $xy$\\6 levels};
  \node[mdsbox, minimum width=30mm, minimum height=9mm] (n2) at (58, -1) {kd-tree\\sliding midpoint, 48 levels};
  \node[mdsbox, minimum width=30mm, fill=white] (nl) at (58, -10) {leaf 128};
  \draw[mdsflow] (n1) -- (n2);
  \draw[mdsflow] (n2) -- (nl);
  \node[mdslabel, anchor=west] at (74, 9) {\schema{qt6\_kd48l128sm}};
  % wide top
  \node[mdslabel] at (108, 16) {wide top};
  \node[mdsbox, minimum width=30mm] (w1) at (108, 9) {STR, 64 children\\1 level};
  \node[mdsbox, minimum width=30mm, minimum height=9mm] (w2) at (108, -1) {octree, points' centre\\24 levels};
  \node[mdsbox, minimum width=30mm, fill=white] (wl) at (108, -10) {leaf 512};
  \draw[mdsflow] (w1) -- (w2);
  \draw[mdsflow] (w2) -- (wl);
  \node[mdslabel, anchor=west] at (124, 9) {\schema{str64\_ot24l512tc}};
  \node[mdslabel, anchor=west, text width=150mm] at (-2, -18)
    {the node at depth $d$ is split by the block whose levels cover $d$; the last block is open-ended (48 or 24 levels, more than any cloud needs); one leaf capacity per schema};

  % ------------------------------------------------------------ right: the pool
  \begin{scope}[shift={(0, -40)}]
    \node[mdslabel, anchor=west] at (-2, 13) {(b) one pool for every schema: 32-byte hot nodes in 64-byte sibling groups, depth first; cold data apart; points in tree order};
    \node[mdslabel, anchor=west] at (-2, 8.5) {hot node: tight box (6 floats) $\cdot$ first child or first point $\cdot$ count $|$ internal flag};
    % hot node lines
    \foreach \i in {0,...,9} {
      \draw[mdsline, fill=mdsfill] (\i*14, 0) rectangle ++(13, 5);
      \draw[mdsline] (\i*14+6.5, 0) -- ++(0, 5);
    }
    % root group (slot 0 pads, slot 1 root) and three sibling groups
    \draw[mdsaccent, thick] (0, -0.8) rectangle ++(13, 6.6);
    \node[mdslabel, text=mdsaccent] at (6.5, -3) {root};
    \draw[mdsaccent, thick] (14, -0.8) rectangle ++(27, 6.6);
    \node[mdslabel, text=mdsaccent] at (27.5, -3) {the root's children (4)};
    \draw[mdsaccent, thick] (42, -0.8) rectangle ++(13, 6.6);
    \node[mdslabel, text=mdsaccent] at (48.5, -6) {children of child 1};
    \draw[mdsaccent, thick] (56, -0.8) rectangle ++(13, 6.6);
    \node[mdslabel, text=mdsaccent] at (62.5, -3) {its subtree, depth first};
    \node[mdslabel, anchor=west] at (72, -3) {\ldots then the subtree of child 2, and so on};
    % cold array
    \node[mdslabel, anchor=west] at (-2, -10.5) {cold node (never read by a query): split cell, subtree extent, depth, block};
    \foreach \i in {0,...,9} {
      \draw[mdsline, fill=white] (\i*14, -16) rectangle ++(13, 3.2);
    }
    % point arrays
    \foreach \j/\lab in {0/$x$, 1/$y$, 2/$z$, 3/id} {
      \draw[mdsline, fill=white] (0, -22-\j*3.6) rectangle ++(139, 2.8);
      \node[mdslabel, anchor=east] at (-0.5, -20.6-\j*3.6) {\lab};
      \draw[mdsout, fill=mdsoutfill] (42, -22-\j*3.6) rectangle ++(27, 2.8);
    }
    \node[mdslabel, anchor=west, text=mdsout] at (70, -36.5) {the points of one subtree are one contiguous run of every array: a leaf scan streams};
  \end{scope}
\end{tikzpicture}%
}
\caption{(a) The schema grammar. A schema is a sequence of blocks, each a primitive with a split rule and a number of levels, and one leaf capacity; a single structure is a schema with one block. (b) Every schema compiles into the same pool: 32-byte hot nodes holding the tight box and the child or point range, two per cache line, children contiguous as a sibling group that starts on a line boundary, subtrees laid out depth first; the data a query never reads lives in a parallel cold array; the points are reordered into tree order as separate coordinate arrays, so every subtree is a contiguous run. Composition changes the split rule at some depths and nothing else.}
\label{fig:grammar}
\end{figure*}


Table~\ref{tab:designspace} lists what an index variant may change. Most of it is familiar: the primitive, its split rule, the leaf capacity and the depth. One option needs a short explanation, because hybrid indexes have been proposed for point clouds several times~\cite{wang2021hybridpointcloud,park2021nestedoctree,ogayar2023nestedlidar}: \emph{nesting}, the composition of several primitives in one hierarchy.

\begin{table}[t]
\centering\footnotesize
\caption{What an index variant may select. A variant is one primitive or a sequence of primitives (nesting), one leaf capacity, a node order and, when out of core, a placement and a cache size.}
\label{tab:designspace}
\begin{tabular}{p{0.24\linewidth}p{0.68\linewidth}}
\toprule
Choice & Options \\
\midrule
Primitive & kd-tree; quadtree; octree; regular grid ($3^3$ cells); hierarchical grid ($4^3$ cells); sort-tile-recursive tree (8, 27 or 64 children) \\
Split rule & kd-tree: median of the longest axis, sliding midpoint, cell centre; quadtree and octree: cell centre or centre of the points' box \\
Leaf capacity & 16 to 32\,768 points \\
Depth & a level count per primitive; the last primitive is open-ended (deep enough for any cloud here) \\
Nesting & a sequence of primitives, each over a number of levels (e.g.\ a 6-level quadtree over a kd-tree) \\
Node order & depth-first or page-blocked \\
Placement & in memory; leaves on a device; hierarchy and leaves on a device; LRU cache size \\
Memory pages & 4\,KiB or 2\,MB (huge pages); child prefetch on or off \\
\bottomrule
\end{tabular}
\end{table}

A nested variant is a sequence of blocks, each a primitive with a split rule and a number of levels, sharing one leaf capacity: the node at depth $d$ is split by the block whose levels cover $d$ (Fig.~\ref{fig:grammar}a). A single structure is the one-block case. In this engine a nested variant is not a different storage scheme: it has the same hot nodes, the same sibling groups, the same leaf arrays and the same traversal as a single structure, and the only thing a second block changes is the split rule, hence the fan-out, at some depths. Two conventions keep the comparison with single structures fair: every variant has one leaf capacity, so a nested variant cannot win by giving its last block a different leaf; and the last block is open-ended, so depth never ends a tree before its leaf capacity does. Without the second convention, nested variants appeared to win on dense scans in our own earlier experiments, by lending a capped kd-tree the levels it was denied; Section~\ref{sec:composition} reports composition under the fair conventions.

\subsection{The search as a source of the best structure}\label{sec:search}

For each cloud and workload, a search on one half of the query trace finds the fastest single structure and the fastest nested variant under the same budget; the other half is held out for every number in this paper. It combines a fixed grid of 55 tuned single structures (every primitive and split rule at leaf capacities from 32 to 32\,768) with a genetic search over one to three blocks, and confirms its finalists by re-timing them in three processes with rotated order. A nested variant is selected only if it is at least 3\% faster than the best single in the median over processes and faster in every process. The search is an instrument here, not a contribution: Section~\ref{sec:rule} shows that a one-feature rule picks nearly as well.

\paragraph{What the engine does not do.} Indexes are static; a changed cloud is rebuilt. Coordinates are single-precision floats in a frame local to the cloud. Reads from the device are synchronous at queue depth one, without prefetching or batching. Points and nodes are not compressed.
